\documentclass[aspectratio=169]{beamer}
\usepackage[utf8]{inputenc}
\usepackage[russian]{babel}
\usepackage[T2A]{fontenc}
\usepackage{tempora} % Отличный Times-подобный шрифт с поддержкой кириллицы
\usepackage{amsmath, amssymb}
\usepackage{graphicx}
\graphicspath{{images/}}


% Настройки темы
\usetheme{Madrid}
\usecolortheme{whale}

% Правила оформления согласно ТЗ
\setbeamerfont{title}{size=\huge, series=\bfseries, family=\rmfamily}
\setbeamerfont{frametitle}{size=\Large, series=\bfseries, family=\rmfamily}
\setbeamerfont{normal text}{size=\large, family=\rmfamily}
\usefonttheme{serif}

\title[Решение задач релаксации]{Разработка программных солверов на суперкомпьютерных системах для
решения кинетического уравнения в задачах релаксации}
\author[Володин М.В.]{Докладчик: Володин Максим \\
Б02-206к \\
Физтех-школа физики и исследований имени Ландау}
\institute[МФТИ]{Московский физико-технический институт}
\date{02.IV.2026}

\begin{document}

% Слайд 1: Титульный
    \begin{frame}
        \titlepage
    \end{frame}

% Слайд 2: Постановка проблемы
    \begin{frame}{Постановка проблемы}
        \begin{itemize}
            \item \textbf{Объект исследования:} уравнение Больцмана для функции распределения~~$f(t, x, \dot{x})$
            \item \textbf{Сложность:} нелинейный интеграл столкновений высокой кратности.
            \item \textbf{Актуальность:} необходимость точного моделирования разреженных газов в гиперзвуковой
            аэродинамике и вакуумной технике.
        \end{itemize}
    \end{frame}
    
    \begin{frame}{Обзор известных подходов}
        \textbf{Проблема:} нарушение законов сохранения при использовании стандартных методов
        \begin{itemize}
            \item DSMC (\textbf{метод Бёрда}) страдает от статистического шума вблизи равновесия
            \item Метод дискретных скоростей (\textbf{Бродвелл, Кабанн}) разрешает столкновения, идеально попадающие
            в узлы сетки (большинство реальных столкновений запрещено).
            \item Методы с интерполяцией (\textbf{Аристов}): разрешаются все столкновения и интерполируют функцию.
            Не сохраняют энергию и импульс (появляются \("\)паразитные\("\) источники).
        \end{itemize}
    
    \end{frame}

% Слайд 3: Новизна и актуальность
    \begin{frame}{Новизна: консервативный проекционный метод}
        \begin{itemize}
            \item Впервые предложен \textbf{Ф. Г. Черемисиным}.
            \item \textbf{Эффективность:} отказ от статистических методов в пользу строгого счёта.
            \item \textbf{Ключевая идея:} интеграл столкновений представляется как проекция функционала на узлы
            скоростной сетки.
            \item \textbf{Преимущество:} строгое выполнение законов сохранения на уровне разностной схемы.
        \end{itemize}
    \end{frame}

% Слайд 4: Цели и задачи
    \begin{frame}{Цели и задачи работы}
        \textbf{Цель:} моделирование процессов релаксации однокомпонентного газа и смеси методом Черемисина.
        \vfill
        \textbf{Задачи:}
        \begin{enumerate}
            \item Реализация общей численной схемы для задачи релаксации.
            \item Расчёт эволюции анизотропного распределения (модель \textbf{квази-шок}).
            \item Анализ выравнивания температур в смеси «горячего» и «холодного» газов.
            \item Контроль макропараметров и H-функции Больцмана (\textbf{второе начало}) в частности.
        \end{enumerate}
    \end{frame}
    
    \begin{frame}{Обезразмеривание и сетка скоростей}
        Равновесное максвелловское распределение:
        \[\tilde{f_M} = n_i \left(\frac{m_i}{2\pi k_B T_i} \right)^{3/2}~\exp \left(-\frac{m_i v^2}{2 k_B T_i}\right) \]
        
        Переход к безразмерным переменным:
        \begin{gather*}
            \xi_{sq} = \sqrt{\frac{k T_0}{m}}, \quad \lambda = \frac{1}{\sqrt{2}\pi n_0 d^2}, \quad f =
            \frac{\tilde{f}}{n_0 \xi_{sq}^3} \\
            \tau_\lambda = \frac{\lambda}{\xi_{sq}}, \quad t = \frac{\tilde{t}}{\tau_\lambda}, \quad \xi =
            \frac{v}{\xi_{sq}}
        \end{gather*}
        
        В пространстве скоростей вводится сетка, ограниченная \textbf{эллипсоидом обрезания} $\xi_{cut} = 4.8$:
        \begin{gather*}
            \xi_{\alpha_i} = -\xi_{cut} + (\alpha_i - 0.5)\Delta \xi_i, \quad \alpha_i = 1, \dots, N_{\xi_i} \\
            \Delta \xi_i = \frac{2\xi_{cut}}{N_{\xi_i}}, \quad \Delta V = \Delta \xi_x \Delta \xi_y \Delta \xi_z \\
        \end{gather*}
    \end{frame}

% Слайд 5: Математическая модель
    \begin{frame}{Интеграл столкновений}
        Уравнение релаксации
        \[ \frac{df}{dt} = I(f_*, f) \]
        Используется модель \textbf{твёрдых сфер} ($b_{\max} = d$). \\
        Интеграл столкновений
        \[ I(\xi) = \int_0^d \int_0^{2 \pi} \int_{-\xi_{cut}}^{+\xi_{cut}} (f' f'_* - f f_*) g b \, db \, d\epsilon \, d\xi_* \]
    \end{frame}

% Слайд 6: Метод решения: Проекционная схема
    \begin{frame}{Дискретизация и проекция}
        Схема непрерывного счёта (аналог метода \textbf{Гаусса-Зейделя}) с увеличенным шагом:
        \[ f_\gamma^{n+1} = f_\gamma^n + t \delta I_\gamma^{(n)} \]
        Вклад каждого столкновения в точке $\gamma$ проецируется на ближайшие узлы $\lambda, \mu$ с весом:
        \begin{multline*}
            \delta I_\gamma^{(n)} = \frac{b^2_{\max} V_{sph}}{4\sqrt{2}} \frac{N_0}{N_\nu} \Big[ \delta_{\alpha_{\nu,
            \gamma}} + \delta_{\beta_{\nu, \gamma}} - (1 - r_\nu) \big( \delta_{\lambda_{\nu, \gamma}} +
            \delta_{\mu_{\nu, \gamma}} \big)- r_\nu \big( \delta_{\lambda_\nu + s_\nu, \gamma} + \delta_{\mu_\nu -
            s_\nu, \gamma} \big) \Big] \Omega_\nu
        \end{multline*}
        
        Разность \textbf{обратного и прямого} потоков:
        \begin{equation*}
            \Omega_\nu = \Big[ \big(~f_{\lambda_\nu}^{(\nu)} f_{\mu_\nu}^{(\nu)} \big)^{1 - r_\nu}
            \big(~f_{\lambda_\nu + s_\nu}^{(\nu)}~f_{\mu_\nu - s_\nu}^{(\nu)} \big)^{r_\nu} - f_{\alpha_\nu}^{(\nu)}
            f_{\beta_\nu}^{(\nu)} \Big] \Big| \xi_{\alpha_\nu} - \xi_{\beta_\nu} \Big|
        \end{equation*}
    \end{frame}

% Слайд 7: Численное интегрирование: Сетки Коробова
    \begin{frame}{Кубатурные сетки Коробова}
        Для вычисления интеграла используются \textbf{теоретико-числовые сетки Коробова}:
        \[ \mathbf{x}_k = \left\{ \frac{a_1 k}{p}, \frac{a_2 k}{p}, \dots, \frac{a_s k}{p} \right\}, \quad k=1, \dots, p \]
        \begin{itemize}
            \item $p$ — простое число (количество узлов).
            \item Высокая точность для многомерного интегрирования за счёт перетасовки узлов
            \item Погрешность ниже, чем в методе Монте-Карло ($\frac{1}{N}$ против $\frac{1}{\sqrt{N}}$).
        \end{itemize}
    \end{frame}

% Слайд 8: Задача 1: Релаксация анизотропного распределения
    \begin{frame}{Задача 1. Квази-шок}
        \[ f_0 = \frac{0.5}{(2\pi)^{1.5}} \left[ \exp\left(-\frac{(\xi_x - u)^2 + \xi_y^2 + \xi_z^2}{2}\right) +
                                              \exp\left(-\frac{(\xi_x + u)^2 + \xi_y^2 + \xi_z^2}{2}\right) \right] \]
        \begin{itemize}
            \item Параметр $u$.
            \item Физический смысл: два встречных потока газа.
            \item Ожидаемый эффект: перекачка энергии из продольной компоненты тензора температуры $T_{xx}$ в
            поперечные до состояния изотропии.
        \end{itemize}
    \end{frame}

% Слайд 10: Результаты задачи 1
    \begin{frame}{Эволюция температуры задачи 1}
        \begin{figure}
            \includegraphics[width=0.7\textwidth]{t10}
%            \caption{}
            \label{fig:t10}
        \end{figure}
    \end{frame}
    
    \begin{frame}{Эволюция $H$-функции задачи 1}
        \begin{figure}
            \includegraphics[width=0.7\textwidth]{h10}
%            \caption{}
            \label{fig:h10}
        \end{figure}
    \end{frame}

% Слайд 9: Результаты задачи 1: Температуры
    \begin{frame}{Результаты задачи 1}
        \textbf{Температурные зависимости}
        \begin{itemize}
            \item $T_{xx}$~уменьшается; $T_{yy}, T_{zz}$ растут.
            При $t \to \infty$: $T_{xx} = T_{yy} = T_{zz} = T$.
        \end{itemize}
        
        \textbf{Эволюция H-функции Больцмана}
        \[ H = \int_{-\xi_{cut}}^{+\xi_{cut}} f(\xi) \ln f(\xi) \, d\xi \]
        \begin{itemize}
            \item Согласно H-теореме, $dH/dt \le 0$ -- монотонное невозрастание с минимумом в равновесии
            \item Скорость убывания коррелирует с величиной $u$.
        \end{itemize}
    
    \end{frame}

% Слайд 11: Задача 2: Релаксация смеси газов
    \begin{frame}{Задача 2: cмесь «горячего» и «холодного» газов}
        \textbf{Начальное состояние:}
        \[ f_0 = f_M(T=1) + f_M(T=2) \]
        \vfill
        \textbf{Ожидаемые результаты}
        \begin{itemize}
            \item \textit{Парциальные температуры:} $T_{cold}$ и $T_{hot}$ стремятся к общему значению $T_{mix}$.
            \item \textit{Энтропия:} суммарная для системы возрастает.
            \item \textit{Концентрации:} $n_{cold}, n_{hot}$ сохраняются (контроль работы метода).
        \end{itemize}
    \end{frame}
    
    \begin{frame}{Эволюция температур задачи 2}
        \begin{figure}
            \includegraphics[width=0.7\textwidth]{t2}
%            \caption{}
            \label{fig:t2}
        \end{figure}
    \end{frame}
    
    \begin{frame}{Эволюция энтропий задачи 2}
        \begin{figure}
            \includegraphics[width=0.7\textwidth]{s}
%            \caption{}
            \label{fig:s}
        \end{figure}
    \end{frame}
% Слайд 12: Результаты задачи 2: Выравнивание температур
    \begin{frame}{Результаты задачи 2}
        \textbf{Эволюция температур}
        \begin{itemize}
            \item Уменьшение скорости релаксации в ходе эволюции
            \item Итоговая температура $T = \frac{\sum_i n_i T_i}{\sum_i n_i}$.
        \end{itemize}
        
        \textbf{Эволюция энтропий}
        \begin{itemize}
            \item Энтропии холодного газа и смеси растёт
            \item Энтропия горячего газа падает
        \end{itemize}
    \end{frame}

% Слайд 13: Результаты задачи 2: Выравнивание


% Слайд 14: Заключение
    \begin{frame}{Заключение}
        \begin{enumerate}
            \item Успешно применён \textbf{консервативный проекционный метод} для задач релаксации.
            \item Подтверждено \textbf{строгое выполнение} законов сохранения массы, импульса и энергии.
            \item Показано выполнение \textbf{второго начала} в дискретном пространстве скоростей.
            \item Использование \textbf{сеток Коробова} позволило эффективно вычислить интеграл столкновений при
            умеренном числе узлов.
        \end{enumerate}
    \end{frame}

% Слайд 16: Список литературы
    \begin{frame}{Список литературы}
        \begin{thebibliography}{9}
            \bibitem{LL10}
            \textit{Лифшиц Е.М., Питаевский Л. П.} Физическая кинетика М.: Наука.~Гл.~ред.~физ.-мат.~лит., 1979.- 528
            страниц (том X)
            
            \bibitem{dyn}
            \textit{Коган М. Н.} Динамика разреженного газа.~— М.: Наука, 1967
            
            \bibitem{Tcher}
            \textit{Черемисин Ф. Г.} Введение в консервативный проекционный метод вычисления интеграла столкновений и
            метод решения уравнения Больцмана: учебное пособие.~– М.: МФТИ.~– 2015.~– 80 с
        \end{thebibliography}
        \vspace{5mm}
        Работа была выполнена с использованием оборудования центра коллективного пользования «Комплекс моделирования и
        обработки данных исследовательских установок мега-класса» НИЦ «Курчатовский институт», http://ckp.nrcki.ru/
    \end{frame}

% Слайд 17: Спасибо за внимание
    \begin{frame}[plain]
    \centering
    \Huge \textbf{Спасибо за внимание!}
    \end{frame}

\end{document}